\section{Introduction}
\label{sec:introduction}

Modern scientific computing faces an increasingly critical
challenge: the volume of data generated by simulations is growing
faster than the ability to move it efficiently. High-performance
computing applications such as seismic forward modelling, climate
simulation, and molecular dynamics produce datasets ranging from
gigabytes to petabytes, making data movement the dominant
bottleneck in many workflows~\cite{jin2024concealing,lee2002migration}.
GPU-accelerated compression has emerged as an attractive solution,
leveraging the parallel architecture of graphics processors to
achieve throughputs that can exceed storage system
bandwidths~\cite{nvcomp,knorr2021ndzip}.

The high-performance computing landscape, however, has become
increasingly heterogeneous. AMD GPUs now power three of the top ten
systems on the TOP500 list, including the two highest-ranked
exascale systems (El Capitan and Frontier)~\cite{top500nov2025}.
Despite AMD's growing presence in HPC, the availability of optimized
GPU compression tools has not kept pace with the hardware. The
NVIDIA ecosystem benefits from \nvcomp~\cite{nvcomp}, a
vendor-supported, production-ready library available since 2020
that offers architecture-tuned implementations of LZ4, Snappy,
Cascaded, GDeflate, ANS, BitComp, and (since release 5.x) ZFP. In
contrast, the HIP/ROCm ecosystem relies on community-driven efforts
such as \hipcomp{}, which appeared in 2022 and remains experimental,
with limited algorithm coverage, direct-port optimization, and
sparse documentation. This asymmetry creates a tooling vacuum for
the growing number of HPC systems built on AMD GPUs.

A functional hipify-based port of \nvcomp{} establishes an
AMD-compatible baseline~\cite{anon_iccsa2026,kunas2025compressao},
but the resulting implementation is far from optimal: it inherits
launch geometries and host-to-device staging strategies that were
tuned for NVIDIA hardware, leaving substantial throughput on the
table. Recent work on seismic reverse-time migration pipelines
also showed that, even with a correctly ported library, the I/O
subsystem dominates end-to-end time~\cite{kunas2024carla,
boito2017seismic}, which makes the host-to-device staging on the
input path a high-leverage optimization target. Two observations
from a systematic cross-architecture characterization campaign on
AMD's CDNA and RDNA lineups motivate the present work. First, at the kernel level, the default chunk size of 64~KB
inherited from \nvcomp{} leaves the larger AMD parts (notably the
MI300X) wave-count-starved: the number of chunks dispatched is one
to two orders of magnitude smaller than the GPU's wave slot
capacity, and the SIMT pipelines run nearly idle. Second, at the
transfer level, the dominant cost in end-to-end execution is the
host-to-device staging of the input batch, which a naive
per-chunk \texttt{hipMemcpyAsync} from pageable storage caps at
roughly 6~GB/s on PCIe Gen4 -- a fifth of the bus theoretical
peak.

This paper presents \arcto, a profile-driven GPU compression
library for AMD that closes both gaps with two complementary
contributions:

\begin{itemize}
\item A \textbf{wave-saturation chunk-size formula}, derived from
hardware performance counter analysis, that picks the chunk size
which keeps every wave slot busy and unlocks 1.4--2.9x additional
throughput on the compress kernels with no code change other than
a single launch-time parameter.

\item An \textbf{adaptive tiled aggregation API}, derived from a
cost model fitted to the characterization data, that replaces the
single-shot coalesced pinned-host staging used by the published
ICCSA baseline. We show that the single-shot path, when measured
honestly with full lifecycle accounting (\texttt{hipHostMalloc}
plus pageable-to-pinned memcpy plus bulk H2D), is in fact slower
than the scattered-pageable baseline at every input size from
100~MB to 4~GB on RDNA3. The adaptive variant amortizes the
allocation cost across windows of a profile-driven optimal size,
recovering the gain and reaching \textbf{up to 2.26x end-to-end
speedup} over the baseline on MI300X for the 16~GB TTI workload.
\end{itemize}

Both optimizations are derived from performance counter observations
on the AMD lineup, not from arch-agnostic intuition, and both
generalize across the four AMD generations evaluated (MI50, MI210,
MI300X, RX~7900~XT). The cost model behind the adaptive contribution
is also validated by direct measurement: the optimal window size
picked on each architecture by the auto-detection rule matches the
empirically best window size within one octave.

The remainder of the paper is organized as follows.
Section~\ref{sec:background} surveys the AMD GPU landscape, the
batched compression model used by \nvcomp{} and inherited by
\arcto, and the experimental setup used throughout the evaluation.
Section~\ref{sec:arcto-library} \TODO{describes the \arcto{} library
design and its hipify-based porting strategy}.
Section~\ref{sec:optim-chunk-size} \TODO{presents the wave-saturation
chunk-size formula}.
Section~\ref{sec:optim-adaptive-tiled} \TODO{presents the adaptive
tiled aggregation contribution with cost model derivation}.
Section~\ref{sec:evaluation} \TODO{reports the cross-architecture
scaling sweep}.
Section~\ref{sec:related-work} \TODO{positions the work in the
literature}, and Section~\ref{sec:conclusion} concludes.
